% ----------------------------------------------------------------------
%  3.7.4. Тепловое расширение твёрдых тел
% ----------------------------------------------------------------------
\subsection{Тепловое расширение}\label{ssec:expansion}

При нагреве растёт интенсивность колебаний атомов, увеличиваются
межатомные расстояния и размеры тела. Количественно это описывают
коэффициенты линейного ($\alpha$) и объёмного ($\beta$) расширения:
\begin{keyeq}\label{eq:alpha-beta}
\alpha=\frac1l\left(\frac{\partial l}{\partial T}\right)_p,\qquad
\beta=\frac1V\left(\frac{\partial V}{\partial T}\right)_p=\frac1D\left(\frac{\partial p}{\partial T}\right)_V ,
\end{keyeq}
где $D=-V(\partial p/\partial V)_T=1/K$~— модуль всестороннего сжатия. В
изотропных телах (и в кубических кристаллах) расширение одинаково по всем
направлениям, и $\beta=3\alpha$.

\paragraph{Расширение~— следствие ангармонизма}
Будь колебания строго гармоническими, среднее расстояние между атомами не
зависело бы от амплитуды, а значит, и от температуры: гармоническая
потенциальная яма симметрична, и атом одинаково далеко отклоняется в обе
стороны. Реальная кривая взаимодействия \emph{несимметрична}
(рис.~\ref{fig:morse}): отталкивание на малых расстояниях нарастает круто,
а притяжение на больших спадает полого.

При 0~К атомы находятся на расстоянии $r_0$, отвечающем минимуму энергии.
При нагреве атом колеблется, и чем выше температура, тем выше
максимальная потенциальная энергия колебаний ($T_1<T_2<T_3<T_4$ на
рисунке). Каждому уровню энергии отвечают две точки поворота: наибольшего
сближения и наибольшего удаления соседей. Среднее положение атома~—
середина горизонтального отрезка между ними. Из-за асимметрии кривой
середины смещаются вправо: \vazhno{с ростом температуры среднее
межатомное расстояние растёт}.

\begin{figure}[ht]
\centering
\begin{tikzpicture}
\begin{axis}[width=0.72\textwidth, height=6.0cm, xmin=0.6, xmax=3.2, ymin=-1.08, ymax=0.25,
  xlabel={межатомное расстояние $r$}, ylabel={$U/D$},
  xtick={1}, xticklabels={$r_0$}, ytick={-1,0}, axis lines=left,
  legend style={at={(0.98,0.04)}, anchor=south east, font=\scriptsize}]
  \addplot[fmGray, dashed, thin, domain=0.6:3.2, forget plot] {0};
  \addplot[atomLine, very thick] table[x=r, y=U] {figures/th-morse.dat};
  \addlegendentry{реальный потенциал (Морзе)}
  \addplot[fmGreen, thick, dashed, restrict y to domain=-1.1:0.3] table[x=r, y=H] {figures/th-morse.dat};
  \addlegendentry{гармоническое приближение}
  \addplot[fmRed, thick, densely dotted, domain=-0.98:-0.18, samples=40]
    ({1-ln(-x)/3.2}, {x});
  \addlegendentry{средние положения атома}
  % создано tools/gen_thermal.py: уровни энергии и середины отрезков
\draw[fmBlue, thick] (axis cs:0.7954,-0.85) -- (axis cs:1.3062,-0.85);
\fill[fmRed] (axis cs:1.0508,-0.85) circle (1.6pt);
\node[font=\scriptsize, anchor=west] at (axis cs:1.3062,-0.85) {$T_1$};
\draw[fmBlue, thick] (axis cs:0.7095,-0.65) -- (axis cs:1.5597,-0.65);
\fill[fmRed] (axis cs:1.1346,-0.65) circle (1.6pt);
\node[font=\scriptsize, anchor=west] at (axis cs:1.5597,-0.65) {$T_2$};
\draw[fmBlue, thick] (axis cs:0.6532,-0.45) -- (axis cs:1.8458,-0.45);
\fill[fmRed] (axis cs:1.2495,-0.45) circle (1.6pt);
\node[font=\scriptsize, anchor=west] at (axis cs:1.8458,-0.45) {$T_3$};
\draw[fmBlue, thick] (axis cs:0.6101,-0.25) -- (axis cs:2.2563,-0.25);
\fill[fmRed] (axis cs:1.4332,-0.25) circle (1.6pt);
\node[font=\scriptsize, anchor=west] at (axis cs:2.2563,-0.25) {$T_4$};

\end{axis}
\end{tikzpicture}
\caption{Почему тело расширяется при нагреве. Потенциальная энергия пары
атомов (потенциал Морзе, расчёт) и горизонтальные отрезки~— размах
колебаний при температурах $T_1<\ldots<T_4$; красные точки~— середины
отрезков (средние межатомные расстояния). Они смещаются вправо из-за
асимметрии кривой; в симметричной гармонической яме (зелёные штрихи)
середина всегда была бы при $r_0$}
\label{fig:morse}
\end{figure}

\paragraph{Оценка для линейной цепочки}
Разложим силу взаимодействия соседних атомов по смещению $x$ от равновесия,
сохранив квадратичный член:
\[
f(x)=-ax+bx^2,\qquad U(x)=\frac{ax^2}{2}-\frac{bx^3}{3},
\]
где $a$ характеризует жёсткость связи, а $b$~— ангармонизм. В
установившихся колебаниях средняя сила равна нулю:
$\langle f\rangle=-a\langle x\rangle+b\langle x^2\rangle=0$, откуда
\[
\langle x\rangle=\frac ba\langle x^2\rangle .
\]
В главном приближении $\langle U\rangle\approx a\langle x^2\rangle/2$, а
средние кинетическая и потенциальная энергии равны, так что полная
энергия колебаний $E=2\langle U\rangle=a\langle x^2\rangle$. Подставляя
$\langle x^2\rangle=E/a$, получаем смещение и коэффициент расширения
($c_V=\d E/\d T$~— теплоёмкость в расчёте на одну частицу)\ispr{В
оригинале в промежуточных выкладках записано $\langle x^2\rangle\approx\langle
U\rangle/a$ (пропущен множитель 2) и $\langle x^2\rangle\approx2b\langle
U\rangle/a^2$ (следует читать $\langle x\rangle$), а в окончательном выражении
для $\langle x\rangle$ вместо $b$ стоит $g$; итоговая формула для $\alpha$ при
этом верна.}:
\begin{keyeq}\label{eq:alpha-chain}
\langle x\rangle=\frac{b}{a^2}\,E,\qquad
\alpha=\frac1{r_0}\frac{\d\langle x\rangle}{\d T}=\frac{b}{a^2r_0}\,c_V .
\end{keyeq}
\vazhno{Коэффициент расширения пропорционален теплоёмкости} и
ангармонизму $b$. При высоких температурах $c_V=k$, и
$\alpha=bk/(a^2r_0)$; подстановка типичных значений даёт $\alpha$ порядка
$10^{-5}$~К$^{-1}$, что согласуется с опытом. При высоких температурах
$\alpha$ почти не зависит от $T$, а при низких ведёт себя как теплоёмкость
и стремится к нулю при $T\to0$.

\paragraph{Формула Грюнайзена}
Для трёхмерного изотропного тела аналогичное соотношение получил
Грюнайзен:
\begin{keyeq}\label{eq:gruneisen}
\alpha=\frac{\gamma_\text{Г}K}{3V}\,C_V ,
\end{keyeq}
где $K=1/D$~— сжимаемость, $V$~— мольный объём, $C_V$~— мольная
теплоёмкость, $\gamma_\text{Г}$~— \term{постоянная Грюнайзена} (у
большинства металлов 1,5--2,5); она показывает, как частоты колебаний
меняются с объёмом, и не зависит от температуры. Так как $K$ и $V$ тоже
почти не зависят от $T$, \vazhno{температурный ход $\alpha$ повторяет ход
теплоёмкости} (рис.~\ref{fig:cv-debye}): у диэлектриков $\alpha\propto T^3$ при
$T\ll\theta_D$ и $\alpha\approx\text{const}$ при $T>\theta_D$ (если не учитывать
ангармонической и вакансионной добавок к теплоёмкости).

\begin{podrobnee}[Вклад электронов и общая формула]
В металле давление, создаваемое нагревом, складывается из решёточного и
электронного. Для газа свободных электронов $p=2U/3V$, поэтому
$(\partial p/\partial T)_V^\text{эл}=\frac{2}{3V}C_\text{эл}$, и
\[
\alpha=\frac{1}{3DV}\left[\gamma_\text{Г}C_V^\text{реш}+\frac23C_\text{эл}\right].
\]
Поскольку $\gamma_\text{Г}\sim1$, электронный вклад существен только там, где
электронная теплоёмкость сравнима с решёточной,~— при очень низких
температурах ($\lesssim10$~К). В общем случае
\[
\alpha=A\,D\!\left(\frac{\theta_D}{T}\right)+BT+\frac{C}{T^2}\,e^{-E_\text{в}/RT},
\]
где $D(\theta_D/T)$~— функция Дебая (решётка), $BT$~— электроны, последний
член~— равновесные вакансии (заметен только перед плавлением); $A$, $B$,
$C$~— постоянные.
\end{podrobnee}

Пластическая деформация, искажая решётку и усиливая ангармонизм
колебаний, несколько увеличивает коэффициент расширения.

\paragraph{Корреляции с другими свойствами}
Коэффициент расширения определяется межатомным взаимодействием, поэтому
он связан с другими характеристиками связи.
\begin{itemize}
\item \emph{Критерий Линдемана}: кристалл плавится, когда среднее смещение
атомов достигает определённой доли межатомного расстояния. Из него
следует, что объёмное расширение от 0~К до плавления примерно одинаково:
$\beta\Tpl\approx0{,}06$--$0{,}076$ для большинства металлов с кубической и
гексагональной решётками (соответственно $\alpha\Tpl\approx0{,}02$)\ispr{В
оригинале соотношение записано для линейного коэффициента: $\alpha\Tpl=C_0$
с $C_0=0{,}06$--$0{,}076$. Для меди
$\alpha\Tpl=16{,}5\cdot10^{-6}\cdot1358\approx0{,}022$, а
$\beta\Tpl=3\alpha\Tpl\approx0{,}067$, т.\,е. приведённые значения относятся к
объёмному коэффициенту.}.
Тугоплавкие металлы поэтому расширяются слабо, легкоплавкие~— сильно.
\item Коэффициент расширения тем меньше, чем прочнее отдельная связь:
$\alpha\propto z/(2\Delta H_\text{субл})$, где $\Delta H_\text{субл}$~— энергия
сублимации, $z$~— координационное число, так что $2\Delta H_\text{субл}/z$~—
энергия, приходящаяся на одну связь\ispr{В оригинале соотношение записано
в виде $\alpha\sim(2/z)\Delta H_\text{субл}$, т.\,е. с ростом энергии связи
коэффициент расширения рос бы, что противоречит опыту (ср. вольфрам и
свинец).}.
\end{itemize}

\paragraph{Магнитный вклад. Инвар}
У ферро-, ферри- и антиферромагнетиков магнитный порядок даёт большой
вклад в расширение, иногда сравнимый с решёточным. Природа его та же, что
у магнитострикции: обменная энергия зависит от межатомного расстояния,
и знак эффекта определяется знаком производной обменного интеграла по
расстоянию (кривая Бете--Слэтера, рис.~\ref{fig:bethe}).
\begin{itemize}
\item У хрома, марганца, железа (восходящая ветвь кривой) возникновение
спинового порядка \emph{увеличивает} объём. При нагреве к точке Кюри (Нееля)
порядок разрушается, решётка магнитно \emph{сжимается}, и это сжатие
вычитается из обычного теплового расширения. У железа $\alpha$ уменьшается
от $16{,}5\cdot10^{-6}$ до $14{,}7\cdot10^{-6}$~К$^{-1}$ при нагреве от 800 до
1000~К ($T_\text{К}=1043$~К); у хрома $\alpha$ при нагреве к точке Нееля
(${\approx}311$~К) падает почти до нуля.
\item У никеля (нисходящая ветвь) упорядочение \emph{уменьшает} объём, и
вблизи точки Кюри наблюдается заметное увеличение $\alpha$.
\end{itemize}

Особенно силён магнитный вклад в сплавах Fe--(30--40)\,\%\,Ni: он сравним с
решёточным и противоположен ему по знаку. В результате у таких сплавов~—
\term{инваров}~— коэффициент расширения при комнатной температуре близок
к нулю (рис.~\ref{fig:invar}): у классического инвара Fe--36\,\%\,Ni
$\alpha\approx1\cdot10^{-6}$~К$^{-1}$, на порядок меньше, чем у железа и
никеля. Выше точки Кюри (около 500~К) аномалия исчезает, и $\alpha$
приближается к обычному значению ${\approx}18\cdot10^{-6}$~К$^{-1}$. Сплавы
с 34--45\,\%\,Ni ниже 30~К имеют даже \emph{отрицательный} коэффициент
расширения. Инварной аномалией обладают также ферромагнитные сплавы
Fe--Pt, Fe--Pd, Ni--Al, Ni--Pt, антиферромагнитные Fe--Mn, Cr--Co,
некоторые фазы Лавеса (ZrZn$_2$, ZrFe$_2$) и другие сплавы и соединения
переходных металлов.

\begin{primer*}[Зачем нужен инвар]
Открытие инвара (Ш.~Гийом, Нобелевская премия 1920~г.) дало материал для
эталонов длины, мерных лент, маятников точных часов, деталей приборов,
теневых масок кинескопов, мембран резервуаров для сжиженного газа.
Пара «инвар + обычный металл» образует \emph{термобиметалл}, изгибающийся
при нагреве,~— основу термореле и термостатов.
\end{primer*}

\begin{figure}[ht]
\centering
\begin{tikzpicture}
\begin{axis}[width=0.7\textwidth, height=5.6cm, xmin=25, xmax=100, ymin=0, ymax=21,
  xlabel={Ni, \% (по массе)}, ylabel={$\alpha$, $10^{-6}$~К$^{-1}$},
  xtick={30,36,40,50,60,70,80,90,100}]
  \addplot[fmBlue, thick, smooth, tension=0.5] coordinates {
    (26,19) (28,18.4) (30,16) (31.5,12.5) (33,7.5) (34.5,3.2) (36,1.2)
    (37.5,2.2) (39,4.0) (42,6.6) (46,8.6) (50,9.9) (55,11.2) (60,12.1)
    (70,13.1) (80,13.6) (90,13.6) (100,13.3)};
  \addplot[fmGray, densely dotted, domain=25:100] {13.3};
  \node[font=\scriptsize, anchor=west] at (axis cs:37,1.0) {инвар};
  \node[font=\scriptsize, fmGray, anchor=south east] at (axis cs:99,13.6) {никель};
\end{axis}
\end{tikzpicture}
\caption{Коэффициент линейного расширения сплавов Fe--Ni при комнатной
температуре (схема по данным Ш.~Гийома): резкий минимум у инвара
(36\,\% Ni) вызван магнитным сжатием решётки, компенсирующим тепловое
расширение}
\label{fig:invar}
\end{figure}

\begin{itogi}
\begin{itemize}
\item Колебания решётки~— $3N$ нормальных колебаний; закон дисперсии
$\omega(k)$; при нескольких атомах в ячейке~— акустические и оптические
ветви. Квант колебания~— фонон ($\hbar\omega$, $\hbar\vb k$); средняя энергия
колебания $\hbar\omega/(e^{\hbar\omega/kT}-1)$.
\item Теплоёмкость решётки: $3R$ при $T\gg\theta_D$ (Дюлонг и Пти),
$\frac{12\pi^4}5R(T/\theta_D)^3$ при $T\ll\theta_D$ (Дебай); модель Эйнштейна
даёт слишком быстрый, экспоненциальный спад.
\item Электронная теплоёмкость $\gamma T$ мала из-за принципа Паули;
$\gamma\propto$ плотности состояний на $\EF$ (велика у переходных металлов);
$C/T=\gamma+AT^2$ при низких $T$.
\item Фазовые превращения: первого рода~— скрытая теплота, $C_p\to\infty$;
второго рода~— конечный скачок $C_p$ ($\lambda$-аномалии).
\item Теплопроводность $\lambda=\frac13c\upsilon l$; у диэлектриков~— фононы,
максимум при низких $T$ ($T^3$ слева, $1/T$ справа); у металлов~—
электроны, $\lambda/\sigma=LT$ (Видеман--Франц--Лоренц).
\item Тепловое расширение~— следствие ангармонизма; $\alpha\propto C_V$
(Грюнайзен); магнитный вклад может его скомпенсировать (инвар).
\end{itemize}
\end{itogi}

\begin{proverka}
\item Почему длина волны нормального колебания цепочки не может быть
меньше $2a$?
\item Чем акустические колебания отличаются от оптических? Есть ли
оптические ветви у кремния, в ячейке которого два одинаковых атома?
\item Почему модель Эйнштейна даёт слишком быстрый спад теплоёмкости при
низких температурах, а модель Дебая~— правильный?
\item Оцените температуру, ниже которой электронная теплоёмкость меди
больше решёточной ($\gamma=0{,}695$~мДж/(моль$\cdot$К$^2$), $\theta_D=343$~К).
\item Почему у переходных металлов $\gamma$ во много раз больше, чем у меди?
\item Какой род имеет магнитное превращение железа в точке Кюри, а какой~—
превращение $\alpha\to\gamma$? Как это видно на кривой $C_p(T)$?
\item Почему теплопроводность диэлектрика имеет максимум и почему его
положение зависит от размеров образца?
\item Почему теплопроводность константана намного ниже, чем у меди и
никеля?
\item Почему тело с гармоническими колебаниями атомов не расширялось бы
при нагреве?
\end{proverka}
